\section{Integration and Evaluation}
\label{sec:evaluation}

This section describes the integration of \texttt{Rust2Isabelle} into Aeneas and evaluates its translation coverage.
The Aeneas regression suite is used to examine individual Rust features, while the SBPF case study considers larger components adapted from a real-world system.
The evaluation concerns the generation and acceptance of Isabelle theories rather than semantic correctness.

\subsection{Integration and Feature Coverage}

% development is based on Aeneas ../..
We implement \texttt{Rust2Isabelle} by extending Aeneas with an Isabelle/HOL backend.
It reuses Charon's extraction from Rust to LLBC and Aeneas's functional translation from LLBC to the Pure AST.
The backend then translates the dependency-ordered Pure declarations and assembles them into an Isabelle theory importing the \rprelude{}.

The Isabelle backend comprises approximately 2.3~kLoC of OCaml, including extensions to the test harness.
The implementation adds Isabelle-specific translation of type expressions and declarations, name generation, and built-in mappings.
Most of the translation is implemented in the main extraction module, which handles expressions, functions, globals, traits, and the assembly of complete Isabelle theories.
In addition, we define the \rprelude{} in approximately 1.5~kLoC of Isabelle/HOL to provide semantic support for translated Rust primitives and built-ins.
Generated theories and case-study source files are excluded from these figures.

We evaluated the backend using regression tests from Aeneas.
Isabelle successfully processed 23 generated theories, covering the Rust features summarized in \autoref{tab:feature-coverage}.
Since Isabelle/HOL and HOL4 are closely related HOL-based theorem provers, the table also compares the feature support of their Aeneas backends.
The HOL4 results are based on its implementation and previously generated theories because it is not currently enabled in the Aeneas test suite.

\begin{table}[!ht]
  \centering
  \caption{Feature comparison between the Isabelle and HOL4 backends. \checkmark~denotes full support, $\triangle$ partial support, and $\times$ unsupported features.}
  \label{tab:feature-coverage}
  \small
  \begin{tabular}{@{}p{0.58\textwidth}>{\centering\arraybackslash}m{0.16\textwidth}>{\centering\arraybackslash}m{0.16\textwidth}@{}}
    \toprule
    \textbf{Rust feature} &
    \textbf{Isabelle} &
    \textbf{HOL4} \\
    \midrule
    Basic types and type parameters
      & \yes & \yes \\
    Structs, enums, and recursive data types
      & \yes & \yes \\
    Scalar arithmetic and comparisons
      & \yes & \yes \\
    Casts, wrapping, and bitwise operations
      & \yes & \partialyes \\
    Expressions and higher-order functions
      & \yes & \yes \\
    Arrays, slices, and vectors
      & \yes & \partialyes \\
    Shared and non-nested mutable borrows
      & \yes & \yes \\
    Basic trait declarations and implementations
      & \yes & \partialyes \\
    Const-generic records and functions
      & \yes & \no \\
    Non-nested loops
      & \yes & \yes \\
    \midrule
    Singly recursive functions
      & \partialyes & \yes \\
    Mutually recursive functions
      & \no & \yes \\
    Standard-library built-ins
      & \partialyes & \partialyes \\
    Nested borrowing and control flow
      & \no & \partialyes \\
    Floating-point and unsafe operations
      & \no & \no \\
    \bottomrule
  \end{tabular}
\end{table}

% \begin{table}[!ht]
%   \centering
%   \caption{Rust feature coverage in the Aeneas regression suite.}
%   \label{tab:feature-coverage}
%   \small
%   \begin{tabular}{@{}p{0.72\textwidth}p{0.20\textwidth}@{}}
%     \toprule
%     \textbf{Rust feature} & \textbf{Status} \\
%     \midrule
%     Basic types and type parameters
%       & Supported \\
%     Structs, enums, and recursive data types
%       & Supported \\
%     Scalar arithmetic, comparisons, casts, and bitwise operations
%       & Supported \\
%     Expressions, pattern matching, and closures
%       & Supported \\
%     Arrays, slices, and vectors
%       & Supported \\
%     Shared and non-nested mutable borrows
%       & Supported \\
%     Basic trait declarations and implementations
%       & Supported \\
%     Const-generic records and functions
%       & Supported \\
%     Non-nested loops
%       & Supported \\
%     \midrule
%     Singly recursive functions
%       & Partial \\
%     Advanced traits and standard-library built-ins
%       & Partial \\
%     \midrule
%     Mutually recursive functions
%       & Unsupported \\
%     Nested mutable borrows and nested loops
%       & Unsupported \\
%     Floating-point and unsafe operations
%       & Unsupported \\
%     \bottomrule
%   \end{tabular}
% \end{table}

The results show that the Isabelle backend supports the main types, expressions, borrowing constructs, and data structures exercised by the selected tests.
Compared with the existing HOL4 backend, it provides broader support for the current Pure AST, particularly for const generics, trait dictionaries, scalar operations, and array, slice, and vector primitives.
Both backends preserve Aeneas's functional representation of mutable borrows.

The HOL4 backend currently provides stronger support for recursion.
Its result type distinguishes successful execution, failure, and divergence, allowing singly and mutually recursive functions to be represented without first establishing termination.
The Isabelle backend does not yet represent divergence explicitly, retains termination obligations for singly recursive functions, and does not generate mutually recursive function definitions.

% we plan to support the unsupported features in Isabelle/HOL
% semantic preservation => future work

% The remaining Isabelle failures mainly involve unsupported library built-ins and more complex nested borrowing or control-flow structures.
% The evaluation therefore demonstrates translation coverage for the supported fragment, but does not establish semantic preservation.

% core components of solana Blockchain
% x86 encoder(dependency) & just-in-time compiler -- with a slight modifying(reason)
\subsection{Case Study: Solana JIT Compiler}
\label{sec:SBPF-case-study}

% safe-mode 干掉
The SBPF virtual machine is a core component used to execute programs in the Solana blockchain.
It contains an interpreter, a verifier, an x86-64 just-in-time (JIT) compiler, and its supporting instruction encoder~\cite{anza2026SBPF}.
The production JIT uses raw pointers, executable memory, and function pointers, which are outside the safe-Rust fragment supported by Aeneas.
We therefore adapt the x86 encoder and construct a safe-mode model of the JIT.
The encoder preserves the production encoding logic but returns bytes in a \texttt{Vec<u8>} instead of writing through raw pointers.
The JIT model supports selected decoded SBPF instructions and replaces executable memory and raw addresses with owned vectors and integer offsets.
\autoref{tab:SBPF-results} summarizes their translation results.

\begin{table}[!ht]
  \centering
  \caption{Translation results for the adapted SBPF components.}
  \label{tab:SBPF-results}
  \small
  \renewcommand{\arraystretch}{1.1}
  \begin{tabular*}{0.92\linewidth}
    {@{\extracolsep{\fill}}lrrr@{}}
    \toprule
    \textbf{Component} &
    \multicolumn{1}{c}{\textbf{Rust LOC}} &
    \multicolumn{1}{c}{\textbf{Theory LOC}} &
    \multicolumn{1}{c}{\textbf{Abstract decls.}} \\
    \midrule
    x86 encoder & 817  & 673  & 11 \\
    JIT model   & 1008 & 5455 & 24 \\
    \bottomrule
  \end{tabular*}
\end{table}

% main idea: not one-to-one correspondence
% Rust LOC explain the code being translated

% The translated code is not very idiomatic because xxx, Aeneas front xxx.
Both adapted components are successfully processed by Charon and Aeneas and produce Isabelle theories.
The reported line counts are physical source and theory lines, so they do not imply a one-to-one correspondence.
For the x86 encoder, extraction starts from the encoding function and retains only the declarations required by that function.
Source-level comments, test code, and declarations unrelated to this entry point do not appear in the generated theory.
Consequently, its 817 lines of Rust produce a shorter theory of 673 lines.

% put result first: 1000 LOC to 5000 LOC
% then reasons
The JIT entry point requires substantially more supporting declarations.
Its theory includes the decoded SBPF instruction types and constants, functions that translate SBPF instructions into x86 instructions, the byte-encoding functions for these x86 instructions, and jump-relocation logic.
Moreover, Charon and Aeneas transform Rust control flow and fallible operations into explicit functional constructs, including pattern matches, result-valued computations, and separate loop body definitions.
These dependencies and explicit representations expand the 1008-line safe JIT model into 5455 lines of Isabelle.

Despite this successful translation, the generated theories are not assumption-free.
They contain 35 abstract declarations in total, represented by Isabelle \texttt{axiomatization} commands for referenced library operations that lack concrete definitions in the current \rprelude{}.
In addition, the safe model also excludes executable-memory allocation, raw memory access, indirect function invocation, instruction metering, and unsupported SBPF instructions.
The case study therefore demonstrates automatic translation of a substantial safe subset of the SBPF JIT, rather than verification of its complete production implementation.

% \begin{table}[!ht]
%   \centering
%   \caption{Static results for the adapted SBPF case study.}
%   \label{tab:SBPF-results}
%   \small
%   \begin{tabular}{@{}p{0.20\textwidth}rrp{0.39\textwidth}r@{}}
%     \toprule
%     \textbf{Component} &
%     \textbf{Rust LOC} &
%     \textbf{Theory LOC} &
%     \textbf{types funcs loops} &
%     \textbf{Axioms} \\
%     \midrule
%     Safe x86 encoder
%       & 817 & 673
%       & Instruction data types and byte encoding
%       & 11 \\
%     Safe JIT model
%       & 239 & 481
%       & Offsets, anchors, program counters, and jump relocation
%       & 10 \\
%     \bottomrule
%   \end{tabular}
% \end{table}